% Part: first-order-logic
% Chapter: beyond
% Section: higher-order-logic

\documentclass[../../../include/open-logic-section]{subfiles}

\begin{document}

\olfileid{fol}{byd}{hol}

\olsection{高阶逻辑}

从一阶逻辑过渡到二阶逻辑，使我们能够在形式语言中谈论一阶论域中对象组成的集合。为何要止步于此呢？例如，三阶逻辑应当能让我们处理对象集合的集合，甚至既包含对象又包含对象集合的集合。而四阶逻辑则允许我们谈论此类对象的集合。你可能已经猜到了，我们可以任意迭代这一想法。

在实践中，高阶逻辑常常用函数而非关系来形式化。（若在函数与关系之间自然的等同意义下看，这种差异是非本质性的。）给定一些基本的“类”$A$、$B$、$C$、……（我们现在称之为“类型”），我们可以通过规定来生成新的类型：

\begin{quote}
如果 $\sigma$ 和 $\tau$ 是有限类型，那么 $\sigma \to \tau$ 也是。
\end{quote}

可以把类型想象成句法“标签”，用来对我们想要纳入论域的对象进行分类；$\sigma \to \tau$ 描述的是这样一类对象，它们是以 $\sigma$ 类型的对象为输入、以 $\tau$ 类型的对象为输出的函数。
例如，我们可能想要有一个真值类型 $\Omega$（“真”和“假”）以及一个自然数类型 $\Nat$。
在这种情况下，你可以把 $\Nat \to \Omega$ 类型的对象看作一元关系，或者说 $\Nat$ 的子集；$\Nat \to \Nat$ 类型的对象是从自然数到自然数的函数；而 $(\Nat \to \Nat) \to \Nat$ 类型的对象则是“泛函”，即以函数为输入、以数为输出的高类型函数。

与二阶逻辑的情况类似，我们可以把高阶逻辑看作一种多类逻辑，其中为我们想要考虑的每种对象都设置一个类。
但通常更清晰的做法是从头开始定义高类型逻辑的语法。
例如，我们可以归纳地定义有限类型的集合，如下所示：

\begin{enumerate}
\item $\Nat$ 是有限类型。
\item 如果 $\sigma$ 和 $\tau$ 是有限类型，那么 $\sigma \to \tau$ 也是有限类型。
\item 如果 $\sigma$ 和 $\tau$ 是有限类型，那么 $\sigma \times \tau$ 也是有限类型。
\end{enumerate}

直观上，$\Nat$ 表示自然数的类型，$\sigma \to \tau$ 表示从 $\sigma$ 到 $\tau$ 的函数的类型，而 $\sigma \times \tau$ 表示由两个对象（一个来自 $\sigma$，一个来自 $\tau$）构成的有序对的类型。
然后我们可以归纳地定义一个项的集合，如下所示：

\begin{enumerate}
\item 对每个类型 $\sigma$，都有一组类型为 $\sigma$ 的变元 $x$, $y$, $z$, \dots。
\item $\Obj 0$ 是类型为 $\Nat$ 的项。
\item $\Obj S$（后继）是类型为 $\Nat \to \Nat$ 的项。
\item 如果 $s$ 是类型为 $\sigma$ 的项，且 $t$ 是类型为 $\Nat \to (\sigma \to \sigma)$ 的项，则 $\Obj{R}_{st}$ 是类型为 $\Nat \to \sigma$ 的项。
\item 如果 $s$ 是类型为 $\tau \to \sigma$ 的项且 $t$ 是类型为 $\tau$ 的项，则 $s(t)$ 是类型为 $\sigma$ 的项。
\item 如果 $s$ 是类型为 $\sigma$ 的项且 $x$ 是类型为 $\tau$ 的变元，那么 $\lambd[x][s]$ 是类型为 $\tau \to \sigma$ 的项。
\item 如果 $s$ 是类型为 $\sigma$ 的项且 $t$ 是类型为 $\tau$ 的项，那么 $\tuple{s, t}$ 是类型为 $\sigma \times \tau$ 的项。
\item 如果 $s$ 是类型为 $\sigma \times \tau$ 的项，那么 $p_1(s)$ 是类型为 $\sigma$ 的项而 $p_2(s)$ 是类型为 $\tau$ 的项。
\end{enumerate}

直观上说，$\Obj{R}_{st}$ 表示递归定义的函数 \begin{align*}
\Obj{R}_{st}(0) & = s \\
\Obj{R}_{st}(x+1) & = t(x, R_{st}(x)),
\end{align*}
$\tuple{s, t}$ 表示第一分量为~$s$、第二分量为~$t$ 的有序对，而 $p_1(s)$ 和 $p_2(s)$ 表示~$s$ 的第一和第二分量（“投影”）。
最后，$\lambd[x][s]$ 表示由 \[
f(x) = s
\] 定义的函数~$f$；
因此条款 (6) 赋予了我们一种概括原则，使得我们可以利用项来定义函数。
公式则由相同类型项之间的等词陈述 $\eq[s][t]$、通常的命题联结词以及高类型量词构造而成。
然后，我们可以取系统的公理为支配上述所定义项的基本等式，再加上带量词和等词的通常逻辑规则。

如果在有限类型系统中增加一个真值类型 $\Omega$，就必须同时包括支配其使用的公理。
事实上，只要足够巧妙，我们完全可以消除复杂的公式，用类型为 $\Omega$ 的项来替代它们！证明系统也可以随之修改。
这样得到的结果本质上就是 Alonzo Church（阿隆佐·丘奇）在 20 世纪 30 年代提出的\emph{简单类型论}。

与二阶逻辑的情况类似，高类型逻辑也有多种可供选择的语义版本。
在完整的版本中，类型 $\sigma \to \tau$ 的变元遍历从类型~$\sigma$ 的对象到类型~$\tau$ 的对象的\emph{所有}函数构成的集合。
正如你可能预料的那样，这种语义太强了，以至于无法容许一个完备且能行的推导系统。
但是，我们可以考虑一种更弱的语义，在这种语义中，一个结构由每个类型 $\tau$ 的元素集合 $T_\tau$ 构成，并配备适当的应用、投影等操作。
如果细节处理得当，就可以为上文描述的那类推导系统获得完备性定理。

高阶逻辑很有吸引力，因为它提供了一个框架，使得我们能够以一种自然的方式嵌入大量数学：从 $\Nat$ 出发，可以定义实数、连续函数等等。
在直觉主义逻辑的背景下，它也特别有吸引力，因为类型有清晰的“构造性”解释。
事实上，人们可以发展构造性的高类型语义（基于直觉主义而非经典逻辑），它能很好地澄清这些构造性解释，并且在许多方面比经典语义的对应物更有趣。

\end{document}